\documentclass[tikz,border=6pt]{standalone}
\usepackage{fontspec}
\usepackage[mathrm=sym]{unicode-math}
\setmainfont{texgyrepagella-regular.otf}[
    BoldFont = texgyrepagella-bold.otf,
    ItalicFont = texgyrepagella-italic.otf,
    BoldItalicFont = texgyrepagella-bolditalic.otf
]
\setmathfont{texgyrepagella-math.otf}[AutoFakeBold = 2.5]
\usepackage{tikz}
\usetikzlibrary{arrows.meta, decorations.markings}

\begin{document}
\begin{tikzpicture}[>=Stealth, line cap=round, line join=round,
  arrowed/.style={
    decoration={
      markings,
      mark=at position 0.55 with {\arrow{Stealth[length=4.8pt, width=3.3pt]}}
    },
    postaction={decorate}
  },
  arrowedpos/.style={
    decoration={
      markings,
      mark=at position #1 with {\arrow{Stealth[length=4.8pt, width=3.3pt]}}
    },
    postaction={decorate}
  }
]

  % ========================================================
  % ROW 1: Bethe-Goldstone integral equation for G
  % ========================================================
  \begin{scope}[shift={(0, 1.4)}]

    % LHS: Box G
    \begin{scope}[shift={(-4.6, 0)}]
      \draw[thick, line width=1pt, arrowedpos=0.6] (-1.0, 0.45) -- (-0.35, 0.16);
      \draw[thick, line width=1pt, arrowedpos=0.6] (-1.0, -0.45) -- (-0.35, -0.16);
      \draw[thick, line width=1pt, arrowedpos=0.55] (0.35, 0.16) -- (1.0, 0.45);
      \draw[thick, line width=1pt, arrowedpos=0.55] (0.35, -0.16) -- (1.0, -0.45);

      \draw[thick, fill=white] (-0.35, -0.35) rectangle (0.35, 0.35);
      \node[font=\large] at (0, 0) {$G$};
    \end{scope}

    % Equals sign
    \node[font=\large] at (-3.2, 0) {$=$};

    % Term 1: Bare interaction \bar{v}
    \begin{scope}[shift={(-1.9, 0)}]
      \coordinate (V1) at (0, 0);
      \draw[thick, line width=1pt, arrowedpos=0.6] (-0.75, 0.45) -- (V1);
      \draw[thick, line width=1pt, arrowedpos=0.6] (-0.75, -0.45) -- (V1);
      \draw[thick, line width=1pt, arrowedpos=0.55] (V1) -- (0.75, 0.45);
      \draw[thick, line width=1pt, arrowedpos=0.55] (V1) -- (0.75, -0.45);
      \filldraw[black] (V1) circle (2pt);
      \node[font=\normalsize, below=5pt] at (V1) {$\bar{v}$};
    \end{scope}

    % Plus sign
    \node[font=\large] at (-0.6, 0) {$+$};

    % Term 2: \bar{v} + intermediate pair + Box G
    \begin{scope}[shift={(1.0, 0)}]
      \coordinate (V2) at (0, 0);
      % Incoming legs
      \draw[thick, line width=1pt, arrowedpos=0.6] (-0.75, 0.45) -- (V2);
      \draw[thick, line width=1pt, arrowedpos=0.6] (-0.75, -0.45) -- (V2);
      \node[font=\normalsize, below=5pt] at (V2) {$\bar{v}$};

      % Bubble to Box G at x = 1.3
      \coordinate (GBtop) at (1.1, 0.16);
      \coordinate (GBbot) at (1.1, -0.16);

      \draw[thick, line width=1pt, arrowedpos=0.55] (V2) to[out=50, in=155] (GBtop);
      \draw[thick, line width=1pt, arrowedpos=0.55] (V2) to[out=-50, in=-155] (GBbot);

      % Box G
      \draw[thick, fill=white] (1.1, -0.35) rectangle (1.8, 0.35);
      \node[font=\large] at (1.45, 0) {$G$};

      % Outgoing legs from Box G
      \draw[thick, line width=1pt, arrowedpos=0.55] (1.8, 0.16) -- (2.45, 0.45);
      \draw[thick, line width=1pt, arrowedpos=0.55] (1.8, -0.16) -- (2.45, -0.45);

      \filldraw[black] (V2) circle (2pt);
    \end{scope}

  \end{scope}

  % ========================================================
  % ROW 2: Ladder expansion of G in powers of \bar{v}
  % ========================================================
  \begin{scope}[shift={(0, -1.3)}]

    % Equals sign
    \node[font=\large] at (-3.2, 0) {$=$};

    % Term 1: Single vertex
    \begin{scope}[shift={(-1.9, 0)}]
      \coordinate (W1) at (0, 0);
      \draw[thick, line width=1pt, arrowedpos=0.6] (-0.75, 0.45) -- (W1);
      \draw[thick, line width=1pt, arrowedpos=0.6] (-0.75, -0.45) -- (W1);
      \draw[thick, line width=1pt, arrowedpos=0.55] (W1) -- (0.75, 0.45);
      \draw[thick, line width=1pt, arrowedpos=0.55] (W1) -- (0.75, -0.45);
      \filldraw[black] (W1) circle (2pt);
    \end{scope}

    % Plus sign
    \node[font=\large] at (-0.6, 0) {$+$};

    % Term 2: Two vertices with ladder rung
    \begin{scope}[shift={(0.8, 0)}]
      \coordinate (L1) at (-0.55, 0);
      \coordinate (L2) at (0.55, 0);

      % Incoming
      \draw[thick, line width=1pt, arrowedpos=0.6] (-1.25, 0.45) -- (L1);
      \draw[thick, line width=1pt, arrowedpos=0.6] (-1.25, -0.45) -- (L1);

      % Bubble
      \draw[thick, line width=1pt, arrowed] (L1) to[out=55, in=125] (L2);
      \draw[thick, line width=1pt, arrowed] (L1) to[out=-55, in=-125] (L2);

      % Outgoing
      \draw[thick, line width=1pt, arrowedpos=0.55] (L2) -- (1.25, 0.45);
      \draw[thick, line width=1pt, arrowedpos=0.55] (L2) -- (1.25, -0.45);

      \filldraw[black] (L1) circle (2pt);
      \filldraw[black] (L2) circle (2pt);
    \end{scope}

    % Plus sign
    \node[font=\large] at (2.6, 0) {$+$};

    % Term 3: Three vertices (two ladder rungs)
    \begin{scope}[shift={(4.7, 0)}]
      \coordinate (M1) at (-1.0, 0);
      \coordinate (M2) at (0, 0);
      \coordinate (M3) at (1.0, 0);

      % Incoming
      \draw[thick, line width=1pt, arrowedpos=0.6] (-1.65, 0.45) -- (M1);
      \draw[thick, line width=1pt, arrowedpos=0.6] (-1.65, -0.45) -- (M1);

      % Bubble 1
      \draw[thick, line width=1pt, arrowed] (M1) to[out=55, in=125] (M2);
      \draw[thick, line width=1pt, arrowed] (M1) to[out=-55, in=-125] (M2);

      % Bubble 2
      \draw[thick, line width=1pt, arrowed] (M2) to[out=55, in=125] (M3);
      \draw[thick, line width=1pt, arrowed] (M2) to[out=-55, in=-125] (M3);

      % Outgoing
      \draw[thick, line width=1pt, arrowedpos=0.55] (M3) -- (1.65, 0.45);
      \draw[thick, line width=1pt, arrowedpos=0.55] (M3) -- (1.65, -0.45);

      \filldraw[black] (M1) circle (2pt);
      \filldraw[black] (M2) circle (2pt);
      \filldraw[black] (M3) circle (2pt);
    \end{scope}

    % Plus and ellipsis
    \node[font=\large] at (7.0, 0) {$+$};
    \node[font=\large] at (7.9, 0) {$\cdots$};

  \end{scope}

\end{tikzpicture}
\end{document}
